% Options for packages loaded elsewhere
% Options for packages loaded elsewhere
\PassOptionsToPackage{unicode}{hyperref}
\PassOptionsToPackage{hyphens}{url}
\PassOptionsToPackage{dvipsnames,svgnames,x11names}{xcolor}
\PassOptionsToPackage{space}{xeCJK}
%
\documentclass[
]{ctexart}
\usepackage{xcolor}
\usepackage[margin=1in]{geometry}
\usepackage{amsmath,amssymb}
\setcounter{secnumdepth}{5}
\usepackage{iftex}
\ifPDFTeX
  \usepackage[T1]{fontenc}
  \usepackage[utf8]{inputenc}
  \usepackage{textcomp} % provide euro and other symbols
\else % if luatex or xetex
  \usepackage{unicode-math} % this also loads fontspec
  \defaultfontfeatures{Scale=MatchLowercase}
  \defaultfontfeatures[\rmfamily]{Ligatures=TeX,Scale=1}
\fi
\usepackage{lmodern}
\ifPDFTeX\else
  % xetex/luatex font selection
  \ifXeTeX
    \usepackage{xeCJK}
    \setCJKmainfont[]{STSong}
  \fi
  \ifLuaTeX
    \usepackage[]{luatexja-fontspec}
    \setmainjfont[]{STSong}
  \fi
\fi
% Use upquote if available, for straight quotes in verbatim environments
\IfFileExists{upquote.sty}{\usepackage{upquote}}{}
\IfFileExists{microtype.sty}{% use microtype if available
  \usepackage[]{microtype}
  \UseMicrotypeSet[protrusion]{basicmath} % disable protrusion for tt fonts
}{}
\makeatletter
\@ifundefined{KOMAClassName}{% if non-KOMA class
  \IfFileExists{parskip.sty}{%
    \usepackage{parskip}
  }{% else
    \setlength{\parindent}{0pt}
    \setlength{\parskip}{6pt plus 2pt minus 1pt}}
}{% if KOMA class
  \KOMAoptions{parskip=half}}
\makeatother
% Make \paragraph and \subparagraph free-standing
\makeatletter
\ifx\paragraph\undefined\else
  \let\oldparagraph\paragraph
  \renewcommand{\paragraph}{
    \@ifstar
      \xxxParagraphStar
      \xxxParagraphNoStar
  }
  \newcommand{\xxxParagraphStar}[1]{\oldparagraph*{#1}\mbox{}}
  \newcommand{\xxxParagraphNoStar}[1]{\oldparagraph{#1}\mbox{}}
\fi
\ifx\subparagraph\undefined\else
  \let\oldsubparagraph\subparagraph
  \renewcommand{\subparagraph}{
    \@ifstar
      \xxxSubParagraphStar
      \xxxSubParagraphNoStar
  }
  \newcommand{\xxxSubParagraphStar}[1]{\oldsubparagraph*{#1}\mbox{}}
  \newcommand{\xxxSubParagraphNoStar}[1]{\oldsubparagraph{#1}\mbox{}}
\fi
\makeatother


\usepackage{longtable,booktabs,array}
\usepackage{calc} % for calculating minipage widths
% Correct order of tables after \paragraph or \subparagraph
\usepackage{etoolbox}
\makeatletter
\patchcmd\longtable{\par}{\if@noskipsec\mbox{}\fi\par}{}{}
\makeatother
% Allow footnotes in longtable head/foot
\IfFileExists{footnotehyper.sty}{\usepackage{footnotehyper}}{\usepackage{footnote}}
\makesavenoteenv{longtable}
\usepackage{graphicx}
\makeatletter
\newsavebox\pandoc@box
\newcommand*\pandocbounded[1]{% scales image to fit in text height/width
  \sbox\pandoc@box{#1}%
  \Gscale@div\@tempa{\textheight}{\dimexpr\ht\pandoc@box+\dp\pandoc@box\relax}%
  \Gscale@div\@tempb{\linewidth}{\wd\pandoc@box}%
  \ifdim\@tempb\p@<\@tempa\p@\let\@tempa\@tempb\fi% select the smaller of both
  \ifdim\@tempa\p@<\p@\scalebox{\@tempa}{\usebox\pandoc@box}%
  \else\usebox{\pandoc@box}%
  \fi%
}
% Set default figure placement to htbp
\def\fps@figure{htbp}
\makeatother





\setlength{\emergencystretch}{3em} % prevent overfull lines

\providecommand{\tightlist}{%
  \setlength{\itemsep}{0pt}\setlength{\parskip}{0pt}}



 


\usepackage{float}
\usepackage{unicode-math}
\floatplacement{figure}{H}
\makeatletter
\@ifpackageloaded{tcolorbox}{}{\usepackage[skins,breakable]{tcolorbox}}
\@ifpackageloaded{fontawesome5}{}{\usepackage{fontawesome5}}
\definecolor{quarto-callout-color}{HTML}{909090}
\definecolor{quarto-callout-note-color}{HTML}{0758E5}
\definecolor{quarto-callout-important-color}{HTML}{CC1914}
\definecolor{quarto-callout-warning-color}{HTML}{EB9113}
\definecolor{quarto-callout-tip-color}{HTML}{00A047}
\definecolor{quarto-callout-caution-color}{HTML}{FC5300}
\definecolor{quarto-callout-color-frame}{HTML}{acacac}
\definecolor{quarto-callout-note-color-frame}{HTML}{4582ec}
\definecolor{quarto-callout-important-color-frame}{HTML}{d9534f}
\definecolor{quarto-callout-warning-color-frame}{HTML}{f0ad4e}
\definecolor{quarto-callout-tip-color-frame}{HTML}{02b875}
\definecolor{quarto-callout-caution-color-frame}{HTML}{fd7e14}
\makeatother
\makeatletter
\@ifpackageloaded{caption}{}{\usepackage{caption}}
\AtBeginDocument{%
\ifdefined\contentsname
  \renewcommand*\contentsname{Table of contents}
\else
  \newcommand\contentsname{Table of contents}
\fi
\ifdefined\listfigurename
  \renewcommand*\listfigurename{List of Figures}
\else
  \newcommand\listfigurename{List of Figures}
\fi
\ifdefined\listtablename
  \renewcommand*\listtablename{List of Tables}
\else
  \newcommand\listtablename{List of Tables}
\fi
\ifdefined\figurename
  \renewcommand*\figurename{Figure}
\else
  \newcommand\figurename{Figure}
\fi
\ifdefined\tablename
  \renewcommand*\tablename{Table}
\else
  \newcommand\tablename{Table}
\fi
}
\@ifpackageloaded{float}{}{\usepackage{float}}
\floatstyle{ruled}
\@ifundefined{c@chapter}{\newfloat{codelisting}{h}{lop}}{\newfloat{codelisting}{h}{lop}[chapter]}
\floatname{codelisting}{Listing}
\newcommand*\listoflistings{\listof{codelisting}{List of Listings}}
\makeatother
\makeatletter
\makeatother
\makeatletter
\@ifpackageloaded{caption}{}{\usepackage{caption}}
\@ifpackageloaded{subcaption}{}{\usepackage{subcaption}}
\makeatother
\usepackage{bookmark}
\IfFileExists{xurl.sty}{\usepackage{xurl}}{} % add URL line breaks if available
\urlstyle{same}
\hypersetup{
  pdftitle={庞加莱猜想与怪球},
  colorlinks=true,
  linkcolor={blue},
  filecolor={Maroon},
  citecolor={Blue},
  urlcolor={Blue},
  pdfcreator={LaTeX via pandoc}}


\title{庞加莱猜想与怪球}
\author{}
\date{2026-03-28}
\begin{document}
\maketitle

\renewcommand*\contentsname{Table of contents}
{
\hypersetup{linkcolor=}
\setcounter{tocdepth}{3}
\tableofcontents
}

\section{Poincare Conjecture}\label{poincare-conjecture}

\subsection{Topological Poincare
Conjecture}\label{topological-poincare-conjecture}

我们把拓扑庞加莱猜想简记为 (GPC) General Poincare Conjecture.

\begin{itemize}
\tightlist
\item
  GPC:

  \begin{itemize}
  \tightlist
  \item
    任意闭的拓扑流形 \(M^n\), 若 \(M^n\) 是同伦球面(\(M\) 同伦等价于
    \(S^n\)), 则 \(M\) 是拓扑球面(同胚于 \(S^n\)).
  \item
    GPC 是对的.
  \end{itemize}
\item
  相比确定同胚型, 确定同伦型是相对容易的.

  \begin{itemize}
  \tightlist
  \item
    主要工具是Whitehead 定理:

    \begin{itemize}
    \tightlist
    \item
      假设 \(X,Y\) 是两个 CW 复形, 若存在连续映射 \(f: X\to Y\), \(f\)
      诱导同伦群上的同构 (称这样的 \(f\) 为弱同伦等价), 则 \(f\)
      是同伦等价.
    \item
      同调版本: 对于单连通的空间 \(X,Y\), 若 \(f: X\to Y\)
      诱导同调群上的同构, 则 \(f\) 是同伦等价.

      \begin{itemize}
      \tightlist
      \item
        Sketch: 使用 Hurewicz 定理.
      \end{itemize}
    \end{itemize}
  \item
    证明拓扑空间 \(X\) 是同伦球面的过程是标准的:

    \begin{itemize}
    \tightlist
    \item
      若 \(X\) 单连通, 且为同调球面, 则 \(X\) 是同伦球面.

      \begin{itemize}
      \tightlist
      \item
        Sketch: 使用 Hurewicz 定理构造映射 \(f\) + 同调版本的 Whitehead
        定理.
      \end{itemize}
    \item
      推论: 单连通的三维流形是同伦球面.
    \end{itemize}
  \item
    猜想: 同调庞加莱猜想

    \begin{itemize}
    \tightlist
    \item
      同调球面是拓扑球面。
    \item
      猜想是错的, 反例: Poincare 构造了 3 维怪球(它是同调球面,
      但不是同伦球), 它的基本群不平凡 (基本群是 \(\tilde{I}\) , 用
      \(\mathbb{Z}_2\) 对 \(A_5\) 进行中心扩张后得到的完美群, 完美群:
      商掉交换子是零)
    \end{itemize}
  \end{itemize}
\item
  weak GPC: 任意闭的\textbf{光滑}流形 \(M^n\), 若 \(M^n\) 是同伦球面, 则
  \(M^n\) 是拓扑球面.

  \begin{itemize}
  \tightlist
  \item
    Smale: weak GPC 在 \(n\geq5\) 的情形下成立.
  \item
    这个结果弱于GPC, 因为它要求流形是光滑的.
    然而我们关心的多数流形都是光滑的, 故这个结果也足够重要.
  \item
    weak GPC 的证明主要依赖于 h-cobordism 定理.
  \end{itemize}
\item
  handle 分解

  \begin{itemize}
  \tightlist
  \item
    handle 分解是胞腔分解在光滑流形中的对应: 注意到, 一般来说,
    粘接不同的胞腔会改变拓扑空间的维数,
    导致得到的胞腔复形一般来说不是流形.
  \item
    handle 分解的定义:

    \begin{itemize}
    \tightlist
    \item
      main idea: 对于一个 \(n\) 维的闭光滑流形 \(M^n\), 我们希望将 \(M\)
      表示成若干个 (有限多个) \(D^n\) 的组合(这里每个 \(D^n\) 被称为
      handle \(h\), 可以定义它的指标, 指标为 k 的 handle 被记为
      k-handle, 对应于胞腔分解中的一个 k-cell).

      \begin{itemize}
      \tightlist
      \item
        每个组成部分 \(D^n\) 都是平凡的, 我们流形 \(M\)
        的复杂程度完全蕴含在了每个 \(D^n\) 的粘接方法中.
      \item
        在考虑粘接 handle 时, 我们应该想象这个操作在胞腔分解中的对应.
        事实上, handle 分解在同伦上给出了一个胞腔分解.
      \end{itemize}
    \item
      在某个光滑流形 \(X^n\) 上粘接 k-handle:

      \begin{itemize}
      \tightlist
      \item
        为了对应粘接 k-cell，我们将 \(D^n\) (把这个 handle 记为 \(h\))
        看作 \(D^k\times D^{n-k}\)，其中第一个分量 \(D^k\)
        用于粘接：我们将 \(D^k\) 的边界 \(S^{k-1}\) 粘接到初始流形 \(X\)
        的边界 \(\partial X\) 上；而 \(D^{n-k}\)
        则用于``补充''，保证粘接得到的结果还是一个流形。
      \item
        具体来说: 给定一个嵌入映射
        \(\varphi: \partial D^k\times D^{n-k}=S^{k-1} \times D^{n-k} \hookrightarrow \partial X\)，粘接后的流形定义为
        \(X \cup_{\varphi} h := X \sqcup (D^k \times D^{n-k}) / \sim\)，其中等价关系
        \(\sim\) 由 \(\varphi\) 给出。
      \item
        在同伦上看, 这个操作对应于粘接一个 k-cell.
      \end{itemize}
    \item
      定义: 光滑流形的 handle 分解

      \begin{itemize}
      \item
        一个 \(n\) 维光滑流形 \(M\) 的一个 handle 分解是指从空集
        \(\varnothing\) 出发，通过有限次粘接 handle 得到 \(M\)
        的一个表示。更细致地说，一个 handle 分解是一个有限序列
        \(\varnothing = M_{-1} \subset M_0 \subset M_1 \subset \cdots \subset M_m = M\)，其中每个
        \(M_j\) 是通过在 \(M_{j-1}\) 上粘接一个 handle 得到的。
      \item
        例子: \(T^2\) 的 handle 分解

        \begin{figure}[H]

        {\centering \pandocbounded{\includegraphics[keepaspectratio]{images/poincare-t2-handle.png}}

        }

        \caption{T² 的 handle 分解}

        \end{figure}%
      \end{itemize}
    \item
      handle 分解的对偶和消去:

      \begin{itemize}
      \tightlist
      \item
        \textbf{add picture}
      \end{itemize}
    \end{itemize}
  \item
    Morse 函数和 handle 分解的存在性:

    \begin{itemize}
    \tightlist
    \item
      结果: 任意紧的(可能带边界)光滑流形上存在 (有限) handle 分解.
    \item
      论证依赖 Morse 理论:

      \begin{itemize}
      \tightlist
      \item
        Morse 函数: 给定一个光滑流形 \(M\) 上的光滑函数
        \(f: M \to \mathbb{R}\)，若其所有临界点(导数为零的点)都非退化
        (二阶导数矩阵, 即Hessian 矩阵
        \(H_{ij}(p) = \frac{\partial^2 f}{\partial x_i \partial x_j}(p)\)
        的行列式不为零) ，则称 \(f\) 为一个 Morse 函数. 我们可以根据
        Hessian 矩阵的 signature(负惯性指数) 来定义临界点的指标。
      \item
        光滑流形 \(M^n\) 上的任意 Morse 函数给出 \(M^n\) 上的一个 handle
        分解. 且指标为 \(k\) 的临界点对应一个 \(k\) -handle 的粘接。
      \item
        光滑流形 \(M^n\) 上总存在 Morse 函数 (事实上, Morse 函数在
        \(C^\infty\left(M\right)\) 中是稠密的,
        \(C^\infty\left(M\right)\) 定义为 \(M\)
        上所有光滑函数构成的空间)
      \item
        结合以上，handle 分解总存在。
      \end{itemize}
    \end{itemize}
  \item
    例子: \(T^2\) 的 handle 分解.
  \end{itemize}
\item
  h-cobordism 和 h-cobordism 定理 (Smale)

  \begin{itemize}
  \tightlist
  \item
    oriented cobordism:

    \begin{itemize}
    \tightlist
    \item
      我们称两个 \(n\) 维定向流形 \(X_1,X_2\) 是定向配边的,
      如果存在一个定向的 \((n+1)\) 维流形 \(M\), 使得在集合上看, 我们有
      \(\partial M=X_1\sqcup X_2\) (仅仅将这里的流形看成集合),
      如果考虑流形的定向, 则我们有
      \(\partial M = X_1 \sqcup (-X_2)\)，其中 \(-X_2\) 表示流形 \(X_2\)
      带相反定向. 有时, 我们也把 \(\left(M;X_1,X_2\right)\)
      称为一个定向配边.

      \begin{itemize}
      \tightlist
      \item
        如果我们将对 \(X_1,X_2,M\) 的定向要求去掉, 并且只要求
        \(\partial M=X_1\sqcup X_2\) 在集合上成立, 则我们称 \(X_1\) 和
        \(X_2\) 是配边的.
      \end{itemize}
    \item
      例子 (定向配边):
    \end{itemize}
  \item
    h-cobordism:

    \begin{itemize}
    \tightlist
    \item
      考虑定向配边 \(X_1,X_2,M\), 以及嵌入映射
      \(\iota_1: X_1\to M, \iota_2: X_2\to M\). 我们称 \((M;X_1,X_2)\)
      是一个 h-配边，如果每个嵌入 \(\iota_i\) 都是同伦等价。

      \begin{itemize}
      \tightlist
      \item
        这里的 \(h\) 代表同伦。
      \end{itemize}
    \end{itemize}
  \item
    h-cobordism Theorem:

    \begin{itemize}
    \item
      定理: 设 \(\left(M;X_1,X_2\right)\) 是一个紧的、单连通的
      h-配边，且 \(\dim X_1 \geq 5\)，则 \(M\) 微分同胚于一个柱面
      \(X_1 \times [0,1]\)。

      \begin{itemize}
      \tightlist
      \item
        拓扑 h-cobordism Thm: 保持假设不变, 推出 \(M\) 同胚于柱面
        \(X_1\times [0,1]\)
      \end{itemize}
    \item
      推论: 若 \(\left(M;X_1,X_2\right)\) 如上所述, 则 \(X_1\) 与
      \(X_2\) 微分同胚.
    \item
      推论: 对于维数大于等于 5 的情形 \[
        \{  拓扑球的集合 \}/ \text{Diffeo} \cong \{  同伦球的集合 \}/\text{h-cobordism}
        \]

      \begin{itemize}
      \tightlist
      \item
        证明: 利用 GPC+ h-cobordism 定理.
      \end{itemize}
    \item
      定理的证明思路: 我们对 pair \(\left(M,X_1\right)\) 做 handle 分解
      (即将 \(M\) 写成在 \(X_1\times[0,1]\) 上粘接 handles
      的形式)。利用单连通和同伦等价的强限制, 我们可以证明这个 handle
      分解实际上没有多余的 handles, 即 \(M\cong X_1\times[0,1]\).
    \end{itemize}
  \end{itemize}
\item
  历史: GPC 的证明

  \begin{itemize}
  \tightlist
  \item
    weak GPC (\(n\geq5\)): Smale
  \item
    GPC \(\left(n\geq5\right)\): Newman by Connell

    \begin{itemize}
    \tightlist
    \item
      key word: Engulfing method (吞噬法)
    \end{itemize}
  \item
    GPC \(\left(n=4\right)\): Freedman

    \begin{itemize}
    \tightlist
    \item
      证明概要:

      \begin{itemize}
      \tightlist
      \item
        Freedman: 拓扑 h-cobordism 定理在 \(\dim X_1=4\) 的情形下成立.
      \item
        Wall: 设 \(M,N\) 是单连通的闭四维流形, 若它们之间有等价的
        intersection form (一个同伦不变量, 后面介绍), 则它们之间存在一个
        h-配边。
      \item
        推论: 对于单连通闭四维流形 \(M\)，\(M\) 的拓扑完全由它上面的
        intersection form 决定, 从而同伦等价蕴含同胚。这蕴含了 GPC, n=4.
      \end{itemize}
    \end{itemize}
  \item
    GPC \(\left(n=3\right)\): Perelman 通过证明几何化猜想而完成。
  \item
    GPC \(\left(n<3\right)\): trivial.
  \end{itemize}
\end{itemize}

还有很多数学家做出了重要的贡献, 即使我们没有一一列举.

\subsection{Smooth Poincare Conjecture and Exotic
Sphere}\label{smooth-poincare-conjecture-and-exotic-sphere}

\begin{itemize}
\tightlist
\item
  SPC: 光滑庞加莱猜想

  \begin{itemize}
  \tightlist
  \item
    SPC: 设 \(M^n\) 是光滑闭流形, 若它是同伦球面, 则它也是微分球面
    (微分同胚于标准 \(S^n\)).
  \item
    一般的 SPC 是错的:

    \begin{itemize}
    \tightlist
    \item
      Milnor 在 1956 年发现了第一个怪球 \(S^7\).
    \item
      事实上, SPC 在多数维数上失效。
    \item
      既然一般来说,(拓扑) \(S^n\) 上的光滑结构不止一种,
      我们考虑其上所有光滑结构的集合, 记为 \(\theta\left(n\right)\).
    \end{itemize}
  \item
    对于 \(n\geq5\), 我们知道如何计算 \(\theta\left(n\right)\)

    \begin{itemize}
    \tightlist
    \item
      至少, 我们将计算 \(\theta\left(n\right)\)
      归结成了计算球面的稳定同伦群(难算).
    \end{itemize}
  \item
    对于 \(n\leq3\): 流形上的光滑结构唯一, 故 GPC 推出 SPC.
  \item
    对于 \(n=4\): \textbf{我们完全不知道 \(\theta\left(4\right)\), SPC 4
    是开放问题}。
  \end{itemize}
\item
  拥有唯一光滑结构的奇数维拓扑球面，只有 1 维、3 维、5 维以及 \textbf{61
  维}。其余所有奇数维球面都必然存在怪异微分结构. (王国祯, 徐宙利, 2016)
\end{itemize}

下文中, 我们主要讨论 \(\theta\left(n\right),n\ge 5\) 的计算问题.

\begin{itemize}
\tightlist
\item
  定义: Kervaire-Milnor group

  \begin{itemize}
  \tightlist
  \item
    对于 \(n\geq5\) 的情形下, 我们已经得到了双射:
    \(\{  拓扑球的集合 \}/ \text{Diffeo} \cong \{  同伦球的集合 \}/\text{h-cobordism}\).
    这其实是个群同构. 左边的群, 记为 \(\Theta_n\), 被称为
    Kervaire-Milnor 群.
  \end{itemize}
\end{itemize}

于是, 我们只要关心 \(\Theta_n\) 的计算问题.

\section{\texorpdfstring{Computation of
\(\Theta_n\)}{Computation of \textbackslash Theta\_n}}\label{computation-of-theta_n}

\subsection{\texorpdfstring{\(\Theta_n\) as a
Group}{\textbackslash Theta\_n as a Group}}\label{theta_n-as-a-group}

\begin{itemize}
\tightlist
\item
  定向配边群:

  \begin{itemize}
  \tightlist
  \item
    回顾定向配边, 固定维数 \(n\), 我们考虑所有定向 \(n\)
    维流形的定向配边类构成的集合, 记为 \(\Omega_n\).
  \item
    \(\Omega_n\) 作为群:

    \begin{itemize}
    \tightlist
    \item
      加法: \([\Sigma_1]+[\Sigma_2]:=[\Sigma_1\sqcup\Sigma_2]\).
    \item
      逆元: \(-[\Sigma]:=[-\Sigma]\)
    \item
      零元素: \([\emptyset]\)

      \begin{itemize}
      \tightlist
      \item
        注意 \([\partial M]=[\emptyset]\)
      \end{itemize}
    \item
      容易验证, \(\Omega_n\) 构成 Abel 群.
    \end{itemize}
  \item
    例子: 在 \(n=2\) 维中, \(S^2\times[0,1]\)
  \end{itemize}
\item
  \(\Theta_n\) 作为群

  \begin{itemize}
  \tightlist
  \item
    群结构:

    \begin{itemize}
    \tightlist
    \item
      加法: \([\Sigma_1]+[\Sigma_2]:=[\Sigma_1\# \Sigma_2]\)
    \item
      逆元: \(-[\Sigma]:=[-\Sigma]\)
    \item
      零元: \([S^n]\), 标准光滑球面所在的等价类
    \item
      显然, 加法是交换的. 故为 abel 群.
    \end{itemize}
  \item
    一些说明:

    \begin{itemize}
    \tightlist
    \item
      加法良定义:

      \begin{itemize}
      \item
        lemma: 同伦球的连通和还是同伦球.

        \begin{itemize}
        \tightlist
        \item
          使用命题: 单连通的同调球是同伦球.
        \end{itemize}
      \item
        LEMMA 2.2. Let \(M\), \(M'\) and \(M_r\) be closed and simply
        connected. If \(\Sigma\) is \(h\) -cobordant to \(\Sigma'\) then
        \(\Sigma \# \Sigma_r\) is \(h\) -cobordant to
        \(\Sigma' \# \Sigma_r\).
      \item
        Prove via a picture:

        \pandocbounded{\includegraphics[keepaspectratio]{images/poincare-lemma22.png}}
      \end{itemize}
    \item
      逆元满足 \(-[\Sigma]+[\Sigma]=[S^n]\) :

      \begin{itemize}
      \item
        LEMMA 2.3. A simply connected manifold \(\Sigma\) is \(h\)
        -cobordant to the sphere \(S^n\) if and only if \(\Sigma\)
        bounds a contractible manifold.
      \item
        Sketch:

        \begin{itemize}
        \tightlist
        \item
          假设 \([\Sigma]=[S^n]\), 即 \(\Sigma \sqcup -S^n=\partial M\).
          我们沿着边界上的标准 \(S^n\) 粘回去一个 \(D^{n+1}\), 得到
          \(M'\), 且有 \(\Sigma = \partial M'\)。由于 \(S^n\) 是 \(M\)
          的形变收缩核, 故粘接 disk 之后, 得到的流形 \(M'\) 是可缩的.
        \item
          假设 \(\Sigma\) bounds 一个可缩的 \(M\), 在在 \(M\)
          的内部挖去一个 disk, 我们得到单连通的流形 \(M'\).

          \begin{itemize}
          \tightlist
          \item
            \(S^n\) 是 \(M'\) 的形变收缩核: 考虑 map of pair:
            \(\left(D^{n+1},S^n\right)\hookrightarrow (M,M')\),
            我们可以证明 \(S^n\hookrightarrow M'\) 保持同调群.
          \item
            \(\Sigma\) 是 \(M'\) 的形变收缩核: 对配边
            \(\left(M';\Sigma,S^n\right)\) 使用 Poincare 对偶,
            我们得到相对同调群同构
            \(H_*\left(M',\Sigma\right)\cong H^{n+1-*}\left(M',S^n\right)\)
            (用 handle 分解的对偶性质) , 从而
            \(H_*\left(M',\Sigma\right)\cong0\) , 至此我们说明了
            \(\Sigma\hookrightarrow M'\) 保持同调群, 故它是同伦等价.
          \end{itemize}
        \end{itemize}
      \item
        LEMMA 2.4. If \(\Sigma\) is a homotopy sphere, then
        \(\Sigma \# (-\Sigma)\) bounds a contractible manifold.
      \item
        Sketch: 为了做连通和, 在 \(\Sigma\) 上面挖掉一个 disk.
        这样会得到带边流形 \(\Sigma-D^n\). 我们沿着这个流形的边界旋转
        180 度, 得到流形 \(M\), 如下图所示:

        \pandocbounded{\includegraphics[keepaspectratio]{images/poincare-lemma24.png}}

        这样构造的 \(M\) 显然同伦等价于 \(\Sigma-D^n\)
        (这个同伦可以被造出来), 又注意到, \(\Sigma\) 是同伦球面, 故
        \(\Sigma-D^n\) 可缩, 于是 \(M\) 可缩. 根据 LEMMA 2.3, 我们有
        \([\Sigma]+[-\Sigma]=[S^n]=0\).
      \end{itemize}
    \end{itemize}
  \end{itemize}
\item
  Relation between these two Groups

  \begin{itemize}
  \tightlist
  \item
    \(\Omega^h_n\) 中的加法和单位元的定义在 \(\Omega_n\) 中仍然生效:

    \begin{itemize}
    \tightlist
    \item
      \([\Sigma\# \Sigma']=[\Sigma \sqcup\Sigma']\) in \(\Omega_n\):
      考虑在 \(\left(\Sigma\sqcup\Sigma'\right)\times[0,1]\) 上面粘接
      1-handle, 这个 \(\left(n+1\right)\) 维流形给出了我们想要的
      cobordism.
    \item
      \([S^n]=0\) in \(\Omega_n\): 它是 \(D^{n+1}\) 的边界.
    \end{itemize}
  \item
    但我们在 \(\Omega_n^h\) 中只能用连通和和 \([S^n]\) 定义群结构:

    \begin{itemize}
    \tightlist
    \item
      \(\Omega_n\) 中的运算在 \(\Omega^h_n\) 上不封闭.
    \end{itemize}
  \end{itemize}
\end{itemize}

\subsection{Framed cobordism and Stable homotopy
groups}\label{framed-cobordism-and-stable-homotopy-groups}

\subsubsection{Framed cobordism}\label{framed-cobordism}

\begin{itemize}
\tightlist
\item
  稳定可平行化流形

  \begin{itemize}
  \tightlist
  \item
    可平行化: 称一个流形 \(M\) 是可平行化的, 如果它的切丛 \(TM\)
    是一个平凡向量丛.
  \item
    稳定可平行化: 称一个流形 \(M\) 是稳定可平凡化的,
    如果它的切丛直和某个平凡向量丛后成为平凡丛。
  \item
    \(S^n\) 一般不是可平行化流形, 但是稳定可平凡化流形:

    \begin{itemize}
    \tightlist
    \item
      不可平行化: \(S^n\) 只有在维数 \(n=1,3,7\) 可平行化.\\
    \item
      稳定可平行化: \(TS^n\oplus \epsilon^1\) 是平凡丛.
    \end{itemize}
  \end{itemize}
\item
  Framed manifold

  \begin{itemize}
  \tightlist
  \item
    我们知道, 流形 \(M^n\) 可以被嵌入到维数充分大的欧氏空间
    \(\mathbb{R}^{n+k}\) 中. 如果这个嵌入
    \(i:M^n\hookrightarrow \mathbb{R}^{n+k}\) 的法丛是平凡丛,
    我们可以在法丛上选定一个标架场(平凡化).
    法丛上带有平凡化的流形被称为一个 framed manifold.
  \item
    命题: 任意稳定可平行化流形 \(M\), 存在充分大的 \(k\), 以及嵌入
    \(i:M^n \hookrightarrow\mathbb{R}^{n+k}\) 使得法丛是平凡的.
    即对于这种流形, 我们总可以给它一个 frame, 使其成为 framed manifold.

    \begin{itemize}
    \tightlist
    \item
      Sketch:

      \begin{itemize}
      \item
        法丛关系： 因为 \(M^n\) 嵌在 \(\mathbb{R}^{n+k}\)
        中，它的切丛和法丛的直和一定是整个平坦宇宙的切丛（即平凡丛）：
        \[
          TM \oplus \nu^k \cong \epsilon^{n+k}
          \]
      \item
        稳定可平行化前提： 你假设了 \(M\)
        是稳定可平凡化的。它的代数定义是：只要给切丛加上某个足够大维度
        \(a\) 的平凡丛，它就会变成一个平凡丛： \[
          TM \oplus \epsilon^a \cong \epsilon^{n+a}
          \]
      \item
        我们在第 1 个公式的两边，同时加上一个 \(\epsilon^a\)： \[
          (TM \oplus \epsilon^a) \oplus \nu^k \cong \epsilon^{n+k} \oplus \epsilon^a
          \] 把括号里的东西替换成第 2 个公式的结果： \[
         \epsilon^{n+a} \oplus \nu^k \cong \epsilon^{n+k+a}
         \] 这个公式的白话翻译是： 你的法丛
        \(\nu^k\)，如果再给它强行加上一个极其庞大的平凡向量丛
        \(\epsilon^{n+a}\)，它俩混在一起后，就彻底变成了一个完美的、没有打结的平凡丛！在拓扑学里，这就叫证明了
        \(\nu^k\) 是 ``稳定平凡的 (Stably trivial)''。
      \end{itemize}
    \end{itemize}
  \item
    并不是所有的流形上都可以定义 frame. 球面 \(S^n\) 上可以定义 frame.
    同伦球面上也可以定义 frame.
  \end{itemize}
\item
  Framed cobordism:

  \begin{itemize}
  \tightlist
  \item
    Framed cobordism 是两个 framed manifold
    之间的等价关系，它考虑的是在更高一维的流形中，两个 framed manifold
    是否可以通过一个 framed cobordism 连接起来。我们将这个配边群记为
    \(\Omega^{\text{fr}}_n\).
  \end{itemize}
\end{itemize}

我们从定义中发现, framed manifold \(\left(M,f\right)\) 和嵌入的欧氏空间
\(\mathbb{R}^{n+k}\) 的维数 \(k\) 有关系. 即 framed manifold
是在稳定意义下的 (我们可以在所有的 framed manifold
构成的集合上定义等价关系: 由''稳定''给出,
然后考虑在这个等价关系下的等价类. ) 事实上,
这种稳定的现象还出现在了同伦群的计算中.

\begin{itemize}
\tightlist
\item
  球面的稳定同伦群

  \begin{itemize}
  \tightlist
  \item
    \(\pi_k\left(S^k\right)\cong\mathbb{Z},\forall k\geq1\)

    \begin{itemize}
    \tightlist
    \item
      换一种写法:
      \(\lim_{ k \to \infty }\pi_{0+k}\left(S^k\right)=\mathbb{Z}\)
    \end{itemize}
  \item
    这可以被推广 (by suspension theorem):

    \begin{itemize}
    \tightlist
    \item
      \(\lim_{ k \to \infty }\pi_{n+k}\left(S^k\right)\) 收敛 (当 \(k\)
      充分大时, 这个群不依赖于 \(k\) 的选取)
    \item
      将这个群定义为球面的 \(n\) 阶稳定同伦群, 记为 \(\Pi_n\)。
    \end{itemize}
  \end{itemize}
\item
  定理 (Pontryagin-Thom):

  \begin{itemize}
  \tightlist
  \item
    Framed cobordism 群 \(\Omega_n^{\text{fr}}\) 与球面 \(S^n\)
    的稳定同伦群 \(\Pi_n\) 是同构的.
  \item
    我们给出映射 \(\Omega_n^{\text{fr}}\to \Pi_n\), 它是著名的
    Pontryagin-Thom构造, 在不同的 cobordism 类有众多推广.

    \begin{itemize}
    \tightlist
    \item
      Pontryagin-Thom 构造: 对于
      \(i:M^n \hookrightarrow \mathbb{R}^{n+k}\) 和 framing \(f\),
      粗略地说, 我们将给出从环境流形 \(\mathbb{R}^{n+k}\) 到 \(M^n\)
      的一个坍缩映射(事实上, 我们给出从环境流形 \(\mathbb{R}^{n+k}\)
      的紧致化到 \(M^n\) 的管状邻域的紧致化的坍缩映射).\\
    \item
      \(M\) 的管状邻域的紧致化: 我们考虑 \(M^n\) 在 \(\mathbb{R}^{n+k}\)
      中的管状邻域 \(N\), 它是平凡 disk bundle. 将这个平凡 disk bundle
      沿着边界 \(\partial N=M^n\times S^{k-1}\) 粘起来, 得到
      \(N/\partial N=M^n\times D^k/M^n\times S^{k-1}=M\times S^k/M=M_+\wedge S^k\),
      这个操作其实就是对于 disk bundle \(M\times D^k\) 做紧致化
      \(\left(M\times D^k\right)^+=M^+\wedge S^k=M_+\wedge S^k\).
      这也被称为这个 disk bundle 的 Thom 空间.
    \item
      环境流形 \(\mathbb{R}^{n+k}\) 的紧致化是 \(S^{n+k}\).
    \item
      将紧致化看作函子, 它是反变的, 原来的管状邻域的含入映射
      \(M\times D^k\hookrightarrow \mathbb{R}^{n+k}\) 在这里给出
      \(S^{n+k}\to M_+\wedge S^k=M^n\times S^k/M\). 将这个映射与
      \(M\times S^k/M\to S^k\) 复合, 我们得到映射 \(S^{n+k} \to S^k\),
      从而得到一个元素 \(\alpha \in \pi_{n+k}(S^k)\). 这个元素
      \(\alpha\) 在稳定同伦群 \(\Pi_n\) 中的像就是 Pontryagin-Thom
      构造给出的对应.
    \end{itemize}
  \end{itemize}
\end{itemize}

\subsection{Kervaire-Milnor 's exact
sequence}\label{kervaire-milnor-s-exact-sequence}

我们希望把同伦球面''嵌入''到群 \(\Omega_n^{\text{fr}}\) 中去看, 即
\(i:\Theta_n \rightarrow \Omega_n^{\text{fr}}\) ,
从这个映射导出正合序列, 然后使用正合序列计算 \(\Theta_n\),
这样我们就通过一个''已知''的配边群计算另一个未知的配边群. 但是注意到,
映射 \(i\) 性质较差, 比如, 它显然不是满射: 一个同伦球面上面可以有很多个
framing. 我们考虑将 \(\Theta_n\) 映射到 \(\Omega_n^{\text{fr}}\)
上的一个商集.

\begin{itemize}
\tightlist
\item
  \(\Omega_n^{\text{fr}}\) 上的群作用:

  \begin{itemize}
  \tightlist
  \item
    \(\Omega_n^{\text{fr}}\) 上有一个自然的 \(\pi_n(SO)\) 作用.
  \item
    回顾流形上的一个 framing 就是法丛上的一个平凡化,
    即法丛上的一组标架场. 为了从一个 framing 得到另一个 framing,
    只要旋转这个标架场, 即在每个法丛的纤维上作用一个旋转矩阵. 这对应了
    \(\pi_n\left(SO\right)\) 上的元素.
  \item
    给定一个 framing 作为 base, 我们可以定义群作用
    \(J:\pi_n\left(SO\right)\to \Omega_n^{\text{fr}}\),
    把这个映射的像记为 \(J_n\). 商掉这个群作用得到商空间
    \(\Omega_n^{\text{fr}}/J_n\) , 它只依赖于底流形 \(M\), 而不再依赖
    framing.
  \end{itemize}
\end{itemize}

我们将映射改写成
\(\phi:\Theta_n\to \Omega_n^{\text{fr}}/J_n=\Pi_n/J_n\), 并由此构造
Kervaire-Milnor 正合序列.

\begin{itemize}
\tightlist
\item
  映射 \(\phi\) 的核：

  \begin{itemize}
  \tightlist
  \item
    映射 \(\phi\) 的核由那些同伦球面构成，它们作为带 framing 的流形在
    \(\Omega_n^{\text{fr}}\) 中平凡，即应该是某个 framed manifold
    的边界，或者说某个稳定可平行化流形的边界。
  \item
    LEMMA 3.4: 一个稳定可平行化的流形, 如果是带边的,
    那么它也是可平行化的 (parallelizable).
  \item
    我们将这些 bounded, parallelizable 流形的边界构成的子群记为
    \(\Theta^{bp}_n\), 于是 \(\ker \phi=\Theta^{bp}_n\).
    这个群是容易计算的.
  \end{itemize}
\end{itemize}

于是我们有了正合序列: \[
        0 \to \Theta^{bp}_n \to \Theta_n \xrightarrow{\phi} \Pi_n/J_n \to \text{Coker} \phi \to 0
        \]

\begin{itemize}
\tightlist
\item
  映射 \(\phi\) 的像:

  \begin{itemize}
  \item
    我们有如下的结果：
  \item
    对于 \(n \neq 2 \pmod{4}\)，每个 \(\Pi_n\)
    中的元素都可以被同伦球面实现，即 \(\Theta_n \to \Pi_n/J_n\)
    是满射，于是我们有正合列： \[
      0 \to \Theta_n^{bp} \to \Theta_n \to \Pi_n/J_n \to 0
      \]
  \item
    对于 \(n \equiv 2 \pmod{4}\)（此时 \(\Theta_n^{bp} = 0\)），映射
    \(\Theta_n \to \Pi_n/J_n\) 是否为满射由 Kervaire
    不变量决定，即我们有正合列 \[
      0 \to \Theta_{4k-2} \to \Pi_{4k-2}/J_{4k-2} \to \mathbb{Z}/2
      \] 这里多数情况 \(\Pi_{4k-2}/J_{4k-2} = 0\).
    \(\Pi_{4k-2}/J_{4k-2}\neq 0\) 当且仅当
    \(k = 1,2,4,8,16,32\)，对应的流形维数为 \(2,6,14,30,62,126\)
    (\(126\) by 林伟南, 王国祯, 徐宙利
    \href{https://arxiv.org/abs/2412.10879}{On the Last Kervaire
    Invariant Problem} )
  \end{itemize}
\end{itemize}

\section{\texorpdfstring{Computation of some
\(\Theta_n\)}{Computation of some \textbackslash Theta\_n}}\label{computation-of-some-theta_n}

\subsection{\texorpdfstring{Formulas for \(J_n\) and
\(\Theta_n^{bp}\)}{Formulas for J\_n and \textbackslash Theta\_n\^{}\{bp\}}}\label{formulas-for-j_n-and-theta_nbp}

我们记 \(B_n\) 为第 \(n\) 个伯努利数.

\begin{itemize}
\tightlist
\item
  \(J_n\) 的计算：

  \begin{itemize}
  \item
    \(J_n\) 是循环群, 且阶数满足:

    \[|J_n| =
      \begin{cases}
      \text{denominator}\left( \frac{B_{2k}}{4k} \right) & \text{for } n = 4k - 1, \\
      2 & \text{for } n \equiv 0,1 \pmod{8}, \\
      1 & \text{for } n \equiv 2,4,5,6 \pmod{8}.
      \end{cases}\]
  \end{itemize}
\item
  \(|\Theta_n^{bp}|\) 的计算：

  \begin{itemize}
  \tightlist
  \item
    对于 \(n \geq 4\)，群 \(\Theta^{bp}_n\)
    是有限循环群，且有显式的生成元。
  \item
    具体地：

    \begin{itemize}
    \tightlist
    \item
      当 \(n\) 为偶数时，\(\Theta^{bp}_n\) 是平凡群；
    \item
      当 \(n = 4k - 3\) 时，\(\Theta^{bp}_n\) 是平凡群或 \(2\)
      阶循环群；
    \item
      当 \(n = 4k - 1 > 3\) 时，\(\Theta^{bp}_n\) 是阶为
      \(2^{2k-2}(2^{2k-1} - 1) \cdot \text{numerator}\left(\frac{4B_{2k}}{k}\right)\)
      的循环群。
    \end{itemize}
  \end{itemize}
\end{itemize}

\subsection{Examples}\label{examples}

\begin{itemize}
\item
  例子：计算 \(\Theta_6\)

  \begin{itemize}
  \item
    由于 \(6\equiv 2 \pmod{4}\) 我们有正合列:\\
    \[
      0 \to \Theta_6 \to \Pi_6/J_6 \xrightarrow{\Phi_2} \mathbb{Z}/2
      \]
  \item
    查表得到 \(\Pi_6\cong\mathbb{Z}/2\)
  \item
    根据公式, \(J_6=0\). 于是 \(\Pi_6/J_6 \cong \mathbb{Z}/2\).
  \item
    而 \(\Phi_2\) 将生成元映为 \(1 \in \mathbb{Z}/2\), 故为同构. 于是
    \(\Theta_6\cong 0\), 即 \(S^6\) 上只有一个光滑结构.
  \end{itemize}
\item
  例子: 计算 \(\Theta_7\)

  \begin{itemize}
  \tightlist
  \item
    我们用这个方法计算一下著名的 7 维怪球数量 (\(\Theta_7\))。
  \end{itemize}

  \begin{enumerate}
  \def\labelenumi{\arabic{enumi}.}
  \tightlist
  \item
    检查维度：\(n=7 = 4\times2-1\)，属于第一种情况（\(n \equiv 3 \pmod 4\)），正合列分裂。
  \end{enumerate}

  \[
    |\Theta_7| = |\Theta_7^{bp}| \cdot |\Pi_7 / J_7|
    \]

  \begin{enumerate}
  \def\labelenumi{\arabic{enumi}.}
  \setcounter{enumi}{1}
  \tightlist
  \item
    计算 \(|\Theta_7^{bp}|\)：

    \begin{itemize}
    \item
      代入公式 \(k=2\)：

      \[
        |\Theta_7^{bp}| = 2^{2\times2-2}(2^{2\times2-1}-1) \cdot \text{numerator}\left(\frac{4B_4}{2}\right)
        \]
    \item
      第一部分：\(2^2 (2^3 - 1) = 4 \times 7 = 28\)
    \item
      第二部分：\(B_4 = -1/30\)，所以
      \(\frac{4\times1/30}{2} = \frac{1}{15}\)。分子是 1。
    \item
      结果：\(|\Theta_7^{bp}| = 28 \times 1 = 28\)。
    \end{itemize}
  \item
    \textbf{计算 \(|\Pi_7 / J_7|\)}：

    \begin{itemize}
    \tightlist
    \item
      查阅同伦群表，\(\Pi_7 \cong \mathbb{Z}_{240}\)。
    \item
      \(J_7\) 的阶数由 \(\text{denominator}(B_{2k}/4k)\) 决定（其中
      \(n=4k-1\)，即 \(k=2\)）。\(B_4 = \frac{1}{30}\)，所以
      \(\frac{B_4}{8} = \frac{1}{240}\)，分母是 240，因此
      \(|J_7| = 240\)。由于 \(\Pi_7 \cong \mathbb{Z}_{240}\)，故
      \(\Pi_7 = J_7\)，即 \(\Pi_7/J_7\) 是平凡群。
    \item
      可以看到对于 \(n=7\)，\(|\Theta_7|\) 对应的值就是
      \(b_2 = 28\)。这意味着对于 \(n=7\)，\(\Pi_7/J_7\) 是平凡的（或者说
      \(\Pi_7 = J_7\)）。
    \end{itemize}
  \item
    \textbf{结论}：\(|\Theta_7| = 28 \times 1 = 28\)。这意味着有 28
    个微分同胚类，其中 1 个是标准球面，27 个是怪球。
  \end{enumerate}
\end{itemize}

\section{Further Discussion}\label{further-discussion}

\begin{itemize}
\tightlist
\item
  流形的分类

  \begin{itemize}
  \tightlist
  \item
    按照同伦、同胚、微分同胚分类

    \begin{itemize}
    \tightlist
    \item
      这引出了庞加莱猜想.
    \end{itemize}
  \item
    按照配边分类

    \begin{itemize}
    \tightlist
    \item
      即计算配边群(环)
    \item
      主要结果是 Pontryagin-Thom 定理
    \end{itemize}
  \end{itemize}
\item
  配边与同调
\end{itemize}

\subsection{Cobordism Classification of
Manifolds}\label{cobordism-classification-of-manifolds}

我们可以考虑多种配边: 非定向配边(法丛群 \(O\left(n\right)\)),
定向配边(法丛结构群 \(SO\left(n\right)\)), framed 配边(法丛结构群
\(\{ e \}\)), 复配边(法丛结构群 \(U\left(n\right)\))等.

流形在这些配边下的分类是已知的, 我们有结果:

\begin{itemize}
\tightlist
\item
  \textbf{Theorem 1.5.10} (Pontrjagin, Thom) 设 \(\mathfrak{B}\) 为
  \((B, f)\)-structure, 则存在分次环同构:
  \(\Omega_*^{\mathfrak{B}} \cong \pi_*(M\mathfrak{B})\).

  \begin{itemize}
  \tightlist
  \item
    配边群还有环结构: 乘法定义为流形的笛卡尔积 \(X\times Y\).
  \item
    这里 \(\mathfrak{B}\) 可以取成 \(O,SO,\{e\},U\),
    对应上述不同的配边。
  \item
    同构左侧的环可以被完全计算. 即完全分类了流形的上述配边类.
  \end{itemize}
\end{itemize}

\subsection{Cobordism and Normal
Bundle}\label{cobordism-and-normal-bundle}

为什么配边分类与法丛密切相关? 因为对于 \(\Sigma=\partial M\) 的情形,
\(\Sigma\) 的法丛给出了 \(M\) 的法丛. 我们通过验证映射
\(\alpha:\Omega_n^o\to \pi_n\left(MO\right)\) 良定义给出说明.

\pandocbounded{\includegraphics[keepaspectratio]{images/poincare-cobordism1.png}}
\pandocbounded{\includegraphics[keepaspectratio]{images/poincare-cobordism2.png}}

\subsection{Computation of Cobordism
Ring}\label{computation-of-cobordism-ring}

\begin{itemize}
\tightlist
\item
  结果:

  \begin{itemize}
  \tightlist
  \item
    不定向配边:

    \begin{itemize}
    \tightlist
    \item
      \(\Omega^O_* \cong \mathbb{Z}_2[x_2, x_4, x_5, x_6, \dots]\)
    \item
      下标的取法为: 下标不等于 \(2^k-1\)
    \item
      对应的生成元:
    \end{itemize}
  \item
    定向配边:

    \begin{itemize}
    \tightlist
    \item
      \(\Omega_*^{SO} \otimes \mathbb{Q} \cong \mathbb{Q}[y_4, y_8, y_{12}, \dots]\)
    \item
      生成元为 \(\mathbb{CP}^2\), \(\mathbb{HP}^2\) 等流形的配边类。
    \end{itemize}
  \item
    framed 配边:

    \begin{itemize}
    \tightlist
    \item
      \(\Omega_n^{fr} \cong \pi_n^S\)
    \item
      对应球面的稳定同伦群。
    \end{itemize}
  \item
    复配边:

    \begin{itemize}
    \tightlist
    \item
      \(\Omega_*^U \cong \mathbb{Z}[z_2, z_4, z_6, z_8, \dots]\)
    \end{itemize}
  \end{itemize}
\end{itemize}

证明方法:

\begin{itemize}
\tightlist
\item
  分类空间同调 \(\xrightarrow{\text{Thom 同构}}\) Thom 谱的同调群
  \(\xrightarrow{\text{确定 } \mathcal{A}\text{-模结构}}\) 算出 Adams
  谱序列 \(E_2\) 页 \(\xrightarrow{\text{塌陷 (微分全为0)}}\) 直接得出
  Thom 谱的同伦群
\end{itemize}

\begin{center}\rule{0.5\linewidth}{0.5pt}\end{center}

\begin{tcolorbox}[enhanced jigsaw, bottomtitle=1mm, toptitle=1mm, coltitle=black, rightrule=.15mm, titlerule=0mm, opacitybacktitle=0.6, breakable, toprule=.15mm, colbacktitle=quarto-callout-note-color!10!white, bottomrule=.15mm, leftrule=.75mm, colframe=quarto-callout-note-color-frame, title=\textcolor{quarto-callout-note-color}{\faInfo}\hspace{0.5em}{本讲要点}, colback=white, arc=.35mm, opacityback=0, left=2mm]

\begin{enumerate}
\def\labelenumi{\arabic{enumi}.}
\tightlist
\item
  七维怪球的数量为28，对应28个微分同胚类，其中27个是怪球。
\item
  流形可按同伦、同胚、微分同胚以及配边关系进行分类。
\item
  Pontryagin-Thom定理建立了配边群与Thom谱稳定同伦群之间的同构。
\item
  配边分类与流形的法丛结构密切相关。
\item
  不同结构（如非定向、定向、framed、复配边）的配边环已被计算。
\end{enumerate}

\end{tcolorbox}




\end{document}
